\documentclass[../main.tex]{subfiles}
\begin{document}

\section{Conclusion}
\label{sec:conclusion}

This thesis investigated what fixed-\(k\) Most Influential Sets reveal about concentrated sensitivity in a prespecified regression coefficient. The results show that coefficient-targeted joint influence has a mechanism-specific interpretation. Under contamination, bad leverage produces close alignment between planted observations and the influential sets identified by the oracle-assisted MIS search, while response outliers and good leverage show substantially weaker correspondence between unusualness and coefficient movement. Influence localisation also describes a different statistical property from estimator stability, since deleting the MIS-selected set can improve or worsen coefficient accuracy depending on the contamination mechanism. Model misspecification generates further variation: several departures produce strong concentrated sensitivity and distinct localisation patterns, while the simulated linear-endogeneity design exhibits a fitted-span invariance in which increasing severity leaves the standardised deletion-sensitivity structure unchanged. The empirical audits examine these sensitivity questions in realised data, where both the magnitude of coefficient movement and the composition of influential sets vary substantially across studies and deletion budgets. Taken together, the evidence supports fixed-\(k\) MIS as a coefficient-targeted localisation framework for identifying where finite-sample dependence is concentrated and for assessing how that dependence relates to outlyingness, estimator stability, and model specification.

Future research could characterise the outcomes, covariates, treatment status, geography, and institutional context of influential observations, including their persistence across deletion budgets. Further extensions could address structured deletion constraints and coefficient-targeted sensitivity in instrumental-variable, GMM, and nonlinear estimators. These directions would broaden the scope of MIS without changing its diagnostic purpose. The central conclusion is that regression coefficients can depend strongly on small groups of observations, and locating those groups reveals finite-sample dependence that conventional outlier diagnostics, robust estimation, and specification analysis do not individually capture.

\end{document}